% draw network

\begin{figure}[!htb]\centering

  \begin{circuitikz}[/tikz/circuitikz/bipoles/length=1cm, line width=0.8pt]

    % grid
    \draw[gray!50!white, line width=0.5pt] (0.0,0.0) to [short] (8.0,0.0);
    \draw[gray!50!white, line width=0.5pt] (0.0,1.6) to [short] (8.0,1.6);
    \draw[gray!50!white, line width=0.5pt] (0.0,3.2) to [short] (8.0,3.2);
    \draw[gray!50!white, line width=0.5pt] (0.0,4.8) to [short] (8.0,4.8);
    \draw[gray!50!white, line width=0.5pt] (0.0,6.4) to [short] (8.0,6.4);
    \draw[gray!50!white, line width=0.5pt] (0.0,8.0) to [short] (8.0,8.0);

    \draw[gray!50!white, line width=0.5pt] (0.0,0.0) to [short] (0.0,8.0);
    \draw[gray!50!white, line width=0.5pt] (1.6,0.0) to [short] (1.6,8.0);
    \draw[gray!50!white, line width=0.5pt] (3.2,0.0) to [short] (3.2,8.0);
    \draw[gray!50!white, line width=0.5pt] (4.8,0.0) to [short] (4.8,8.0);
    \draw[gray!50!white, line width=0.5pt] (6.4,0.0) to [short] (6.4,8.0);
    \draw[gray!50!white, line width=0.5pt] (8.0,0.0) to [short] (8.0,8.0);


    % generators
    \draw (3.2,-1.8) to [sV, fill=magenta!70!cyan] (3.2,-1.2);
    \draw [short] (3.2,-1.2) to [short] (3.2,-1.0);
    \draw (-1.8,1.6) to [sV, fill=green!50!white] (-1.2,1.6);
    \draw [short] (-1.2,1.6) to [short] (-1.0,1.6);
    \draw (9.2,1.6) to [sV, fill=cyan!50!white] (9.8,1.6);
    \draw [short] (9.2,1.6) to [short] (9.0,1.6);
    \draw (4.8,3.6) to [sV, fill=black!60!white] (4.8,3.0);
    \draw [short] (4.8,3.6) to [short] (4.8,3.8);
    \draw (8.4,4.8) to [sV, fill=yellow!60!white] (9.0,4.8);
    \draw [short] (8.0,4.8) to [short] (8.4,4.8);
    \draw (1.6,9.2) to [sV, fill=orange!60!white] (1.6,9.8);
    \draw [short] (1.6,9.0) to [short] (1.6,9.2);


    % large buses
    \draw[line width=2.5pt] (2.7,0) to [short] (3.7,0);
    \draw[line width=2.5pt] (0.0,1.1) to [short] (0,2.1);
    \draw[line width=2.5pt] (8.0,1.1) to [short] (8,2.1);
    \draw[line width=2.5pt] (1.1,8) to [short] (2.1,8);
    \draw[line width=2.5pt] (2.7,3.2) to [short] (3.7,3.2);
    \draw[line width=2.5pt] (2.7,4.8) to [short] (3.7,4.8);
    \draw[line width=2.5pt] (4.3,6.4) to [short] (5.3,6.4);
    \draw[line width=2.5pt] (0.0,5.9) to [short] (-0.0,6.9);
    \draw[line width=2.5pt] (8,4.3) to [short] (8,5.3);
    \draw[line width=2.5pt] (4.3,4.8) to [short] (5.3,4.8);

    % small buses
    \draw[line width=2.5pt] (2.9,-1) to [short] (3.5,-1);
    \draw[line width=2.5pt] (-1.0,1.3) to [short] (-1,1.9);
    \draw[line width=2.5pt] (9,1.3) to [short] (9,1.9);
    \draw[line width=2.5pt] (4.5,3.8) to [short] (5.1,3.8);
    \draw[line width=2.5pt] (1.3,9.0) to [short] (1.9,9.0);

    \draw[line width=2.5pt] (1.75,3.2) to [short] (1.75,2.6);
    \draw[line width=2.5pt] (1.75,4.8) to [short] (1.75,4.2);
    \draw[line width=2.5pt] (-0.95,6.7) to [short] (-0.95,6.1);
    \draw[line width=2.5pt] (4.5,7.35) to [short] (5.1,7.35);

    % trafos
    \draw (3.2,-1.0) to [voosource] (3.2,0.0);
    \draw (-1,1.6) to [voosource] (0,1.6);
    \draw (8,1.6) to [voosource] (9,1.6);
    \draw (1.6,8) to [voosource] (1.6,9);
    \draw (4.8,4.8) to [voosource] (4.8,3.8);

    \draw (-1,6.4) to [voosource] (0,6.4);
    \draw (4.8,6.4) to [voosource] (4.8,7.4);
    \draw (2.7,4.5) to [voosource] (1.7,4.5);
    \draw (2.7,2.9) to [voosource] (1.7,2.9);

    % loads
    \draw[-{Triangle[length=5mm, width=2mm]}, draw=blue!60!white, fill=blue!60!white] (-1,6.4) -- (-1.6,6.4);
    \draw[-{Triangle[length=5mm, width=2mm]}, draw=blue!60!white, fill=blue!60!white] (4.8,7.4) -- (4.8,8.0);
    \draw[-{Triangle[length=5mm, width=2mm]}, draw=red!60!white, fill=red!60!white] (1.7,4.5) -- (1.1,4.5);
    \draw[-{Triangle[length=5mm, width=2mm]}, draw=blue!60!white, fill=blue!60!white] (1.7,2.9) -- (1.1,2.9);

    \draw (3.0,4.8) to [short] (3.0,4.5);
    \draw (3.0,4.5) to [short] (2.7,4.5);

    \draw (3.0,3.2) to [short] (3.0,2.9);
    \draw (3.0,2.9) to [short] (2.7,2.9);

    % lines
    \draw (3.15,0) to [short] (3.15,3.2);
    \draw (3.25,0) to [short] (3.25,3.2);
    \draw (3.15,3.2) to [short] (3.15,4.8);
    \draw (3.25,3.2) to [short] (3.25,4.8);
    \draw (4.6,4.8) to [short] (4.6,5.1);
    \draw (4.6,5.1) to [short] (3.4,5.1);
    \draw (3.4,5.1) to [short] (3.4,4.8);
    \draw (3.15,4.8) to [short] (3.15,5.1);
    \draw (3.25,4.8) to [short] (3.25,5.1);
    \draw (3.15,5.1) to [short] (4.55,6.1);
    \draw (3.25,5.1) to [short] (4.65,6.1);
    \draw (4.55,6.1) to [short] (4.55,6.4);
    \draw (4.65,6.1) to [short] (4.65,6.4);
    \draw (5.0,6.4) to [short] (5.0,6.1);
    \draw (5.0,6.1) to [short] (7.7,5.1);
    \draw (7.7,5.1) to [short] (8.0,5.1);
    \draw (0,6.4) to [short] (0.3,6.4);
    \draw (0.3,6.4) to [short] (3.0,5.1);
    \draw (3.0,5.1) to [short] (3.0,4.8);


    % legend
    \draw[-{Triangle[length=5mm, width=2mm]}, draw=blue!60!white, fill=blue!60!white] (3.0,8.4) -- (4.0,8.4);
    \draw[-{Triangle[length=5mm, width=2mm]}, draw=red!60!white, fill=red!60!white] (3.0,9.1) -- (4.0,9.1);
    \draw (3.2,9.8) to [sV, fill=magenta!70!cyan] (3.8,9.8);

    \draw (6.0,9.8) to [sV, fill=black!60!white] (6.6,9.8);
    \draw (6.0,9.1) to [sV, fill=yellow!60!white] (6.6,9.1);
    \draw (6.0,8.4) to [sV, fill=green!50!white] (6.6,8.4);

    \draw (9.0,9.8) to [sV, fill=orange!60!white] (9.6,9.8);
    \draw (9.0,9.1) to [sV, fill=cyan!50!white] (9.6,9.1);

    \node at (4.73, 9.8) {\footnotesize Nuclear};
    \node at (4.66, 9.1) {\footnotesize Load I};
    \node at (4.72, 8.4) {\footnotesize Load II};

    \node at (7.64, 9.8) {\footnotesize Dismantled};
    \node at (7.49, 9.1) {\footnotesize Intercon.};
    \node at (7.28, 8.4) {\footnotesize Wind};

    \node at (10.25, 9.8) {\footnotesize Solar};
    \node at (10.40, 9.1) {\footnotesize Storage};

    \draw[gray!50!white, line width=0.5pt] (8.5,8.0) to [short] (10.1,8.0);
    \draw[gray!50!white, line width=0.5pt] (8.5,6.4) to [short] (10.1,6.4);
    \draw[gray!50!white, line width=0.5pt] (8.5,6.4) to [short] (8.5,8.0);
    \draw[gray!50!white, line width=0.5pt] (10.1,6.4) to [short] (10.1,8.0);

    \node at (9.30, 6.2) {\footnotesize 50 km};
    \node[rotate=90] at (10.3, 7.2) {\footnotesize 50 km};

    \draw [fill=gray, opacity=0.2, line width=0.01pt] (2.95,10.15) rectangle (11.0,8.05);
    \draw [fill=gray, opacity=0.2, line width=0.01pt] (11.0,8.05) rectangle (8.05,6.05);

    % buses nodes and labels
    \node at (3.6, 0.2) {8};
    \node at (3.6, -1.0) {1};
    \node at (3.6, 4.6) {9};
    \node at (1.5, 4.75) {2};
    \node at (3.6, 3.0) {10};
    \node at (1.5, 3.15) {3};
    \node at (4.4, 6.6) {11};
    \node at (4.4, 7.6) {4};
    \node at (0.2, 6.7) {12};
    \node at (-1.2, 6.7) {5};
    \node at (5.2, 4.6) {13};
    \node at (5.3, 3.8) {7};
    \node at (7.8, 4.6) {6}; 



    % \node at (8,8) {\begin{tikzpicture} \draw[black, line width=0.8pt, fill=magenta!70!cyan] circle (4pt); \end{tikzpicture} \footnotesize Nuclear}

    % \node at (8,8) {\begin{tikzpicture} \draw[-{Triangle[length=5mm, width=2mm]}] (8,8.4) -- (9,8.4) ; \end{tikzpicture} \footnotesize Load I}





  \end{circuitikz}

  \caption{Overview of the network}
  \label{fig:net1}
\end{figure}


% voltages
% \begin{table}[!htb]\footnotesize
%     \centering
%     \begin{tabular}{|l|l|l|l|l|l|l|l|l|l|l|l|}
%     \hline
%     \diagbox[width=1.5cm, height=0.7cm]{Hour}{Bus} & 1 & 2 & 3 & 4 & 5 & 6 & 8 & 9 & 10 & 11 & 12 \hline
%         0 & 1.050 & 0.958 & 0.986 & 0.967 & 0.929 & 1.000 & 1.030 & 0.982 & 0.997 & 0.979 & 0.942 \hline
%         1 & 1.050 & 0.971 & 0.995 & 0.978 & 0.945 & 1.000 & 1.034 & 0.992 & 1.006 & 0.989 & 0.957 \hline
%         2 & 1.050 & 0.978 & 1.001 & 0.985 & 0.955 & 1.000 & 1.037 & 0.998 & 1.011 & 0.995 & 0.966 \hline
%         3 & 1.050 & 0.981 & 1.003 & 0.987 & 0.959 & 1.000 & 1.038 & 1.001 & 1.013 & 0.997 & 0.969 \hline
%         4 & 1.050 & 0.982 & 1.003 & 0.988 & 0.960 & 1.000 & 1.038 & 1.001 & 1.013 & 0.998 & 0.970 \hline
%         5 & 1.050 & 0.978 & 1.001 & 0.985 & 0.955 & 1.000 & 1.037 & 0.998 & 1.011 & 0.995 & 0.966 \hline
%         6 & 1.050 & 0.963 & 0.989 & 0.971 & 0.935 & 1.000 & 1.032 & 0.985 & 1.000 & 0.983 & 0.947 \hline
%         7 & 1.050 & 0.932 & 0.965 & 0.944 & 0.894 & 1.000 & 1.021 & 0.960 & 0.979 & 0.958 & 0.909 \hline
%         8 & 1.050 & 0.907 & 0.946 & 0.923 & 0.862 & 1.000 & 1.013 & 0.940 & 0.962 & 0.939 & 0.880 \hline
%         9 & 1.050 & 0.896 & 0.937 & 0.913 & 0.847 & 1.000 & 1.008 & 0.930 & 0.954 & 0.930 & 0.865 \hline
%         10 & 1.050 & 0.891 & 0.934 & 0.909 & 0.840 & 1.000 & 1.007 & 0.926 & 0.950 & 0.926 & 0.859 \hline
%         11 & 1.050 & 0.892 & 0.935 & 0.910 & 0.842 & 1.000 & 1.007 & 0.927 & 0.951 & 0.927 & 0.861 \hline
%         12 & 1.050 & 0.895 & 0.937 & 0.913 & 0.846 & 1.000 & 1.008 & 0.930 & 0.953 & 0.929 & 0.865 \hline
%         13 & 1.050 & 0.897 & 0.939 & 0.914 & 0.849 & 1.000 & 1.009 & 0.931 & 0.955 & 0.931 & 0.867 \hline
%         14 & 1.050 & 0.908 & 0.947 & 0.924 & 0.863 & 1.000 & 1.013 & 0.940 & 0.962 & 0.940 & 0.880 \hline
%         15 & 1.050 & 0.915 & 0.952 & 0.929 & 0.872 & 1.000 & 1.015 & 0.946 & 0.967 & 0.945 & 0.888 \hline
%         16 & 1.050 & 0.918 & 0.955 & 0.933 & 0.877 & 1.000 & 1.016 & 0.949 & 0.969 & 0.948 & 0.893 \hline
%         17 & 1.050 & 0.920 & 0.956 & 0.934 & 0.878 & 1.000 & 1.017 & 0.950 & 0.970 & 0.949 & 0.894 \hline
%         18 & 1.050 & 0.914 & 0.951 & 0.929 & 0.871 & 1.000 & 1.015 & 0.945 & 0.966 & 0.944 & 0.888 \hline
%         19 & 1.050 & 0.894 & 0.936 & 0.912 & 0.845 & 1.000 & 1.008 & 0.929 & 0.953 & 0.929 & 0.863 \hline
%         20 & 1.050 & 0.881 & 0.926 & 0.900 & 0.826 & 1.000 & 1.003 & 0.918 & 0.943 & 0.918 & 0.846 \hline
%         21 & 1.050 & 0.888 & 0.932 & 0.906 & 0.836 & 1.000 & 1.006 & 0.924 & 0.948 & 0.924 & 0.856 \hline
%         22 & 1.050 & 0.914 & 0.952 & 0.929 & 0.871 & 1.000 & 1.015 & 0.945 & 0.966 & 0.944 & 0.888 \hline
%         23 & 1.050 & 0.940 & 0.971 & 0.952 & 0.905 & 1.000 & 1.024 & 0.967 & 0.984 & 0.965 & 0.920 \hline
%     \end{tabular}
%     \caption{Voltage profile, in pu, for 24 hours}
%     \label{tab:volt1}
% \end{table}

\begin{table}[!htb]\footnotesize
    \centering
    \begin{tabular}{|c|ccccccccccc|}
    \hline
    \diagbox[width=1.6cm, height=0.7cm]{\textbf{Hour}}{\textbf{Bus}} & \textbf{1} & \textbf{2} & \textbf{3} & \textbf{4} & \textbf{5} & \textbf{6} & \textbf{8} & \textbf{9} & \textbf{10} & \textbf{11} & \textbf{12} \\ \hline \hline
        \textbf{0} & 1.050 & 0.958 & 0.986 & 0.967 & 0.929 & 1.000 & 1.030 & 0.982 & 0.997 & 0.979 & 0.942 \\
        \textbf{1} & 1.050 & 0.971 & 0.995 & 0.978 & 0.945 & 1.000 & 1.034 & 0.992 & 1.006 & 0.989 & 0.957 \\
        \textbf{2} & 1.050 & 0.978 & 1.001 & 0.985 & 0.955 & 1.000 & 1.037 & 0.998 & 1.011 & 0.995 & 0.966 \\
        \textbf{3} & 1.050 & 0.981 & 1.003 & 0.987 & 0.959 & 1.000 & 1.038 & 1.001 & 1.013 & 0.997 & 0.969 \\
        \textbf{4} & 1.050 & 0.982 & 1.003 & 0.988 & 0.960 & 1.000 & 1.038 & 1.001 & 1.013 & 0.998 & 0.970 \\
        \textbf{5} & 1.050 & 0.978 & 1.001 & 0.985 & 0.955 & 1.000 & 1.037 & 0.998 & 1.011 & 0.995 & 0.966 \\
        \textbf{6} & 1.050 & 0.963 & 0.989 & 0.971 & 0.935 & 1.000 & 1.032 & 0.985 & 1.000 & 0.983 & 0.947 \\
        \textbf{7} & 1.050 & 0.932 & 0.965 & 0.944 & 0.894 & 1.000 & 1.021 & 0.960 & 0.979 & 0.958 & 0.909 \\
        \textbf{8} & 1.050 & 0.907 & 0.946 & 0.923 & 0.862 & 1.000 & 1.013 & 0.940 & 0.962 & 0.939 & 0.880 \\
        \textbf{9} & 1.050 & 0.896 & 0.937 & 0.913 & 0.847 & 1.000 & 1.008 & 0.930 & 0.954 & 0.930 & 0.865 \\
        \textbf{10} & 1.050 & 0.891 & 0.934 & 0.909 & 0.840 & 1.000 & 1.007 & 0.926 & 0.950 & 0.926 & 0.859 \\
        \textbf{11} & 1.050 & 0.892 & 0.935 & 0.910 & 0.842 & 1.000 & 1.007 & 0.927 & 0.951 & 0.927 & 0.861 \\
        \textbf{12} & 1.050 & 0.895 & 0.937 & 0.913 & 0.846 & 1.000 & 1.008 & 0.930 & 0.953 & 0.929 & 0.865 \\
        \textbf{13} & 1.050 & 0.897 & 0.939 & 0.914 & 0.849 & 1.000 & 1.009 & 0.931 & 0.955 & 0.931 & 0.867 \\
        \textbf{14} & 1.050 & 0.908 & 0.947 & 0.924 & 0.863 & 1.000 & 1.013 & 0.940 & 0.962 & 0.940 & 0.880 \\
        \textbf{15} & 1.050 & 0.915 & 0.952 & 0.929 & 0.872 & 1.000 & 1.015 & 0.946 & 0.967 & 0.945 & 0.888 \\
        \textbf{16} & 1.050 & 0.918 & 0.955 & 0.933 & 0.877 & 1.000 & 1.016 & 0.949 & 0.969 & 0.948 & 0.893 \\
        \textbf{17} & 1.050 & 0.920 & 0.956 & 0.934 & 0.878 & 1.000 & 1.017 & 0.950 & 0.970 & 0.949 & 0.894 \\
        \textbf{18} & 1.050 & 0.914 & 0.951 & 0.929 & 0.871 & 1.000 & 1.015 & 0.945 & 0.966 & 0.944 & 0.888 \\
        \textbf{19} & 1.050 & 0.894 & 0.936 & 0.912 & 0.845 & 1.000 & 1.008 & 0.929 & 0.953 & 0.929 & 0.863 \\
        \textbf{20} & 1.050 & 0.881 & 0.926 & 0.900 & 0.826 & 1.000 & 1.003 & 0.918 & 0.943 & 0.918 & 0.846 \\
        \textbf{21} & 1.050 & 0.888 & 0.932 & 0.906 & 0.836 & 1.000 & 1.006 & 0.924 & 0.948 & 0.924 & 0.856 \\
        \textbf{22} & 1.050 & 0.914 & 0.952 & 0.929 & 0.871 & 1.000 & 1.015 & 0.945 & 0.966 & 0.944 & 0.888 \\
        \textbf{23} & 1.050 & 0.940 & 0.971 & 0.952 & 0.905 & 1.000 & 1.024 & 0.967 & 0.984 & 0.965 & 0.920 \\ \hline
    \end{tabular}
    \caption{Voltage profile, in pu, for 24 hours}
    \label{tab:volt1}
\end{table}


\begin{figure}[!htb]\centering
\begin{tikzpicture}
    \begin{axis}[xlabel={Hour}, ylabel={$|V|$ (p.u.)}, grid=both, grid style={line width=.1pt, draw=gray!10}, major grid style={line width=.2pt,draw=gray!50}, xtick distance = 2, ytick distance = 0.02, width=12cm, height=8cm, every plot/.append style={very thick}, xmin = 1, xmax = 24, ymin=0.8, very thick, grid=both, grid style={line width=.4pt, draw=gray!10}, major grid style={line width=.8pt,draw=gray!50}, legend style={at={(1.03,0.14)},anchor=south west}]
        \addplot[color=magenta] table[col sep=semicolon, x=x, y=y1] {Data/phase1/ph1_vpu.csv};
        \addplot[color=red] table[col sep=semicolon, x=x, y=y2] {Data/phase1/ph1_vpu.csv};
        \addplot[color=green] table[col sep=semicolon, x=x, y=y3] {Data/phase1/ph1_vpu.csv};
        \addplot[color=cyan] table[col sep=semicolon, x=x, y=y4] {Data/phase1/ph1_vpu.csv};
        \addplot[color=blue] table[col sep=semicolon, x=x, y=y5] {Data/phase1/ph1_vpu.csv};
        \addplot[color=orange, dashed] table[col sep=semicolon, x=x, y=y6] {Data/phase1/ph1_vpu.csv};
        \addplot[color=magenta, dashed] table[col sep=semicolon, x=x, y=y8] {Data/phase1/ph1_vpu.csv};
        \addplot[color=red, dashed] table[col sep=semicolon, x=x, y=y9] {Data/phase1/ph1_vpu.csv};
        \addplot[color=green, dashed] table[col sep=semicolon, x=x, y=y10] {Data/phase1/ph1_vpu.csv};
        \addplot[color=cyan, dashed] table[col sep=semicolon, x=x, y=y11] {Data/phase1/ph1_vpu.csv};
        \addplot[color=blue, dashed] table[col sep=semicolon, x=x, y=y12] {Data/phase1/ph1_vpu.csv};
\legend{Bus 1, Bus 2, Bus 3, Bus 4, Bus 5, Bus 6, Bus 8, Bus 9, Bus 10, Bus 11, Bus 12};
\end{axis}
\end{tikzpicture}
    \caption{Voltage profile during 24 hours for the initial grid. The low-voltage buses are plotted in solid lines; the high-voltage ones are in dashed lines.}
    \label{fig:vpu_ph1}
  \end{figure}





% line loading
\begin{table}[!htb]\footnotesize
    \centering
    \begin{tabular}{|c|ccccc|}
    \hline
       \diagbox[width=1.6cm, height=0.75cm]{\textbf{Hour}}{\textbf{Load}} & \textbf{8-10} & \textbf{10-9} & \textbf{9-11} & \textbf{9-12} & \textbf{11-6} \\ \hline \hline
        \textbf{0} & 26.455 & 19.983 & 9.608 & 31.041 & 64.359 \\
        \textbf{1} & 24.390 & 18.423 & 9.046 & 28.520 & 59.872 \\
        \textbf{2} & 23.081 & 17.437 & 8.701 & 26.936 & 57.069 \\
        \textbf{3} & 22.536 & 17.027 & 8.559 & 26.280 & 55.912 \\
        \textbf{4} & 22.445 & 16.958 & 8.536 & 26.171 & 55.719 \\
        \textbf{5} & 23.074 & 17.431 & 8.699 & 26.928 & 57.053 \\
        \textbf{6} & 25.751 & 19.450 & 9.415 & 30.178 & 62.821 \\
        \textbf{7} & 30.722 & 23.232 & 10.829 & 36.361 & 73.837 \\
        \textbf{8} & 34.303 & 25.990 & 11.904 & 40.967 & 82.183 \\
        \textbf{9} & 35.951 & 27.271 & 12.414 & 43.140 & 86.055 \\
        \textbf{10} & 36.582 & 27.764 & 12.611 & 43.983 & 87.541 \\
        \textbf{11} & 36.445 & 27.657 & 12.568 & 43.800 & 87.220 \\
        \textbf{12} & 36.002 & 27.311 & 12.430 & 43.209 & 86.176 \\
        \textbf{13} & 35.732 & 27.101 & 12.346 & 42.850 & 85.540 \\
        \textbf{14} & 34.196 & 25.907 & 11.872 & 40.827 & 81.932 \\
        \textbf{15} & 33.261 & 25.184 & 11.587 & 39.611 & 79.743 \\
        \textbf{16} & 32.725 & 24.771 & 11.425 & 38.918 & 78.490 \\
        \textbf{17} & 32.543 & 24.631 & 11.370 & 38.685 & 78.066 \\
        \textbf{18} & 33.360 & 25.261 & 11.617 & 39.739 & 79.975 \\
        \textbf{19} & 36.132 & 27.413 & 12.470 & 43.382 & 86.482 \\
        \textbf{20} & 37.988 & 28.867 & 13.054 & 45.886 & 90.862 \\
        \textbf{21} & 36.992 & 28.085 & 12.739 & 44.534 & 88.508 \\
        \textbf{22} & 33.316 & 25.226 & 11.604 & 39.681 & 79.870 \\
        \textbf{23} & 29.399 & 22.221 & 10.443 & 34.693 & 70.872 \\ \hline
    \end{tabular}
    \caption{Percentual loading of the lines for a full day operation}
\end{table}

\begin{figure}[!htb]\centering
\begin{tikzpicture}
    \begin{axis}[xlabel={Hour}, ylabel={Load (\%)}, grid=both, grid style={line width=.1pt, draw=gray!10}, major grid style={line width=.2pt,draw=gray!50}, xtick distance = 2, ytick distance = 10, width=12cm, height=8cm, every plot/.append style={very thick}, xmin = 1, xmax = 24, ymin=0.0, very thick, grid=both, grid style={line width=.4pt, draw=gray!10}, major grid style={line width=.8pt,draw=gray!50}, legend style={at={(1.03,0.14)},anchor=south west}]
        \addplot[color=magenta] table[col sep=semicolon, x=x, y=l1] {Data/phase1/ph1_line.csv};
        \addplot[color=red] table[col sep=semicolon, x=x, y=l2] {Data/phase1/ph1_line.csv};
        \addplot[color=green] table[col sep=semicolon, x=x, y=l3] {Data/phase1/ph1_line.csv};
        \addplot[color=cyan] table[col sep=semicolon, x=x, y=l4] {Data/phase1/ph1_line.csv};
        \addplot[color=blue] table[col sep=semicolon, x=x, y=l5] {Data/phase1/ph1_line.csv};
\legend{Line 8-10, Line 10-9, Line 9-11, Line 9-12, Line 11-6};
\end{axis}
\end{tikzpicture}
    \caption{Representation of the percentual loading of the lines during 24 hours}
    \label{fig:lines_ph1}
  \end{figure}

% figure with loadings in red intense depending in magnitude?

\clearpage
\newpage
% estimate operating costs
Regarding the operating costs, some estimations are made in order to asses the influence of importing energy and the impact of faults on lines and transformers. 

First, the cost of importing energy depends on the time zone: valley, flat or peak. The analysis that follows considers a working day, which is precisely the date for which the voltages and loading profiles have been shown in Figures \ref{fig:vpu_ph1} and \ref{fig:lines_ph1} respectively. The cost of importing the energy is mathematically expressed as:
\begin{equation}
  C_{imp} = \sum_{k=1}^{n=24}P_{s,k} c(k),
  \label{eq:cimp}
\end{equation}
where $C_{imp}$ stands for the cost of importing energy for a full day, $k$ denotes the index of a given hour, $n$ the total number of hours in a day, $P_{s,k}$ the energy provided by the slack bus (interconnection point) in MWh at hour $k$, and $c(k)$ the cost at a certain hour in \texteuro/MWh. This last term is equal to 45~\texteuro/MWh from 0 to 8 hours, 65~\texteuro/MWh from 8 to 10, 14 to 18 and 22 to 24 hours, and 90~\texteuro/MWh from 10 to 14 and 18 to 22 hours. 

Equation \ref{eq:cimp} can be treated as a weighting sum. With the generation data obtained from the timeseries power flow, the total importing cost of importing energy becomes 418753.03~\texteuro/day, or about 152.84~M\texteuro \ in a full year. It is important to note that the study related to the cost of importing energy is decoupled from the fault analysis. This is not a hundred percent realistic, because it could be that a switch trips and hence a line or a transformer are disconnected. Then, it could happen that the interconnection has to provide more power. However, since the probabilities are extremely low, they are discarded when computing this cost. 

On the other hand, there are the costs due to faults in transformers or lines. About 0.05~failures per km and year are expected in lines, while transformers are meant to fail 0.15~times a year. The penalty for not providing energy is 180 \texteuro/MWh. Given that the length of the lines has an impact on its probability of failure, Table \ref{tab:ll1} shows the length and the subsequent failures per year.

\begin{table}[!htb]\centering
  \begin{tabular}[]{ccc}
    \hline 
    \textbf{Line} & \textbf{Length (km)} & \textbf{Failures/year} \\
    \hline
    8-10 & 100.00 & 5.00 \\
    10-9 & 50.00 & 2.50 \\
    9-11 & 70.71 & 3.54 \\
    9-12 & 111.80 & 5.59 \\
    11-6 & 111.80 & 5.59 \\
    \hline
  \end{tabular}
  \caption{Length and failures per year of all active lines}
  \label{tab:ll1}
\end{table}
Figure \ref{fig:lines_ph1} shows that the line connected to the interconnection point operates at a high load. It is critical to note that if line 8-10 fails, the slack should provide all power, but this would result in exceeding the thermal capacity of the line. Thus, if line 8-10 fails, no power can reach the loads. 
 
% TODO: check if is 2 or 2.5 hours for lines!!!!

In the case of line failures, there is a total disconnection time of 2.5 hours; instead, for transformers it is 8 hours. The expected time that an element will be disconnected in a year is found by multiplying the aforementioned disconnection time by the number of failures that take place during a year. Table \ref{tab:lt1} displays the yearly disconnection time and explains the consequences spotted by running the power flow. This will allow to estimate the penalties due disconnection. 

\begin{table}[!htb]\centering
  \begin{tabular}[]{ccc}
    \hline 
    \textbf{Element} & \textbf{Disconnection time (h)} & \textbf{Consequences} \\
    \hline
    Line 8-10 & 12.50 & No load served - divergence \\
    Line 9-10 & 6.25 & No load served - divergence \\
    Line 9-11 & 8.85 & Loads at buses 2, 3 and 5 unserved \\
    Line 9-12 & 13.98 & Load at bus 5 unserved \\
    Line 11-6 & 13.98 & No load served \\
    Trafo 1-8 & 1.20 & No load served - divergence \\
    Trafo 2-9 & 1.20 & Load at bus 2 unserved \\
    Trafo 3-10 & 1.20 & Load at bus 3 unserved \\
    Trafo 4-11 & 1.20 & Load at bus 4 unserved \\
    Trafo 5-12 & 1.20 & Load at bus 5 unserved \\
    \hline
  \end{tabular}
  \caption{Disconnection time and consequences of losing each element}
  \label{tab:lt1}
\end{table}
Once the unserved loads and the associated disconnection times are known, the next step has to do with applying the penalty as follows:
\begin{equation}
  C_{discon} \approx \sum_{i=1}^{10} \overline{P}_{uns,i} t_{discon,i} C_p,
  \label{eq:discon}
\end{equation}
where $C_{discon}$ is the total disconnection cost, $i$ represents the index of the line or transformer with a total of 10 elements prone to be disconnected (see Table \ref{tab:lt1}), $\overline{P}_{uns,i}$ is the mean unserved power, $t_{discon,i}$ the disconnection time, and $C_p$ the penalty cost to apply. Equation \ref{eq:discon} is an approximation in the sense that the unserved power varies according to the time of the day. To not overcomplicate the problem, it has been decided to pick a representative value such as the average. 

The application of Equation \ref{eq:discon} yields a total yearly penalty cost of 4.99 M\texteuro. For the most part, it is due to the disconnection of lines. Meshing more the system would decrease this cost, but on the other side, it would increase the investment cost. Hence, there is a trade-off between cost and reliability.

All the calculations related to costs have not set an inferior limit to the voltages. However, some of them are likely to be unacceptable in reality. The project will proceed to discuss solutions to this issue in the following phases. 